% Montage d'un titrage avec repérage colorimétrique de l'équivalence (L'essentiel, chapitre 5)
\begin{tikzpicture}[x=1cm,y=1cm,scale=0.92,every node/.style={transform shape=false}]
\pic at (0,0) {potence={hauteur=8.2}};
\pic at (0,0.16) {agitateur={largeur=2.4}};
\filldraw[fill=white,draw=otInk!45,line width=0.4pt] (-1.1,0.66) rectangle (1.1,0.71);
\pic at (0,0.71) {erlen={niveau=0.75,liquide=green!6}};
\pic at (-0.2,0.77) {barreau};
\pic at (0,2.95) {burette={hauteur=5.0,niveau=4.3,liquide=red!45!blue!45}};
\pic at (0,7.6) {pince};
\draw[pcRepere] (0.14,6.6) -- (1.2,7.0) node[pcLegende,anchor=west]{burette graduée :\\solution titrante\\(concentration $c_2$ connue)};
\draw[pcRepere] (0.45,3.67) -- (1.3,4.55) node[pcLegende,anchor=west]{robinet};
\draw[pcRepere] (-0.55,1.1) -- (-1.75,1.75) node[pcLegende,anchor=east]{solution titrée :\\volume $V_1$ prélevé\\à la pipette jaugée};
\draw[pcRepere] (0.5,1.0) -- (1.6,1.9) node[pcLegende,anchor=west]{barreau aimanté};
\draw[pcRepere] (1.0,0.69) -- (1.6,1.2) node[pcLegende,anchor=west]{feuille blanche};
\draw[pcRepere] (-1.0,0.4) -- (-1.75,0.55) node[pcLegende,anchor=east]{agitateur magnétique};
\end{tikzpicture}
\par\smallskip\legende{On verse la solution titrante goutte à goutte, sous agitation, jusqu'au changement de teinte qui signale l'équivalence.}
